\startSection
	\selectlanguage{italian}
	\sectionTitle{L'Inno Nazionale Italiano}
	
	\begin{question}{L'inno nazionale italiano è conosciuto anche con il titolo \enquote{Fratelli d'Italia}. Chi fu l'autore del testo, scritto in un periodo di forte sentimento patriottico durante il Risorgimento?}
		\choice{Giuseppe Verdi}
		\choice{Alessandro Manzoni}
		\choice[correct]{Goffredo Mameli}
		\choice{Giuseppe Garibaldi}
		\choice{Gabriele D'Annunzio}
	\end{question}

	
	\begin{question}{La melodia dell'inno nazionale italiano è stata composta da un musicista ligure, su richiesta dell'autore del testo. Chi fu il compositore della musica dell'Inno di Mameli?}
		\choice{Gioachino Rossini}
		\choice[correct]{Michele Novaro}
		\choice{Giuseppe Verdi}
		\choice{Vincenzo Bellini}
		\choice{Pietro Mascagni}
	\end{question}

	
	\begin{question}{Il testo dell'inno nazionale italiano fu scritto in un momento cruciale del Risorgimento. In quale anno Goffredo Mameli compose le parole de \enquote{Il Canto degli Italiani}?}
		\choice{1846}
		\choice[correct]{1847}
		\choice{1848}
		\choice{1861}
		\choice{1870}
	\end{question}

	
	\begin{question}{L'inno nazionale italiano inizia con un verso che è diventato simbolo dell'identità nazionale. Quale tra i seguenti versi apre ufficialmente l'inno?}
		\choice{\enquote{Italia, Italia, sacra terra}}
		\choice{\enquote{Viva l'Italia unita e forte}}
		\choice[correct]{\enquote{Fratelli d'Italia, l'Italia s'è desta}}
		\choice{\enquote{Gloria all'Italia nostra}}
		\choice{\enquote{Sorgi Italia, dal tuo sonno}}
	\end{question}

	
	\begin{question}{Sebbene fosse già ampiamente utilizzato, \enquote{Il Canto degli Italiani} è stato riconosciuto ufficialmente come inno nazionale solo in epoca repubblicana. In quale anno è avvenuta questa ufficializzazione?}
		\choice[correct]{1946}
		\choice{1948}
		\choice{1861}
		\choice{1922}
		\choice{1870}
	\end{question}
\renderSection
